\section*{Five Open Questions Raised by This Replication}

\begin{enumerate}

\item \textbf{Time-dependent Hamiltonians.} Does the classical brickwall optimization extend
cleanly to $H(t)$, where $U_{\text{target}} = \mathcal{T}\exp(-i\int H(s)\,ds)$ rather than a
fixed $\exp(-itH)$? The paper restricts benchmarks to static $H$; real chemistry / gauge-theory
simulation frequently needs time-dependent driving.
\emph{Next steps:} Extend \texttt{src/replicate.py} to accept $H(t)$ (e.g.\ $g(t) = 1 + 0.5\sin(\omega t)$),
build $U_{\text{target}}$ by Magnus expansion or fine subdivision, rerun at $n=3$, $L \in \{1,2,3\}$,
several $\omega$. Compare $\varepsilon_{\text{approx}}$ vs Trotter II at matched depth. Compute: uicgpu A100 (free).

\item \textbf{Composition with higher-order product formulas.} Can classical brickwall optimization
be composed with Suzuki-fourth-order / Yoshida-sixth-order Trotter to give a hybrid Pareto-better
formula? The paper compares against Trotter I/II/IV separately but never warm-starts the L-BFGS/Newton
optimizer from a Trotter-IV parameter vector or nests optimized short-slices inside an outer Suzuki
composition.
\emph{Next steps:} Add Suzuki-4 to the code; warm-start from its $\theta$ and measure wall-time /
$\varepsilon$ vs random restart; build outer-Suzuki-of-inner-optimized composites at longer $t$
(1.6, 3.2) and compare against single-slice optimized brickwall of matched gate count. Compute: uicgpu.

\item \textbf{Real chemistry / lattice-gauge benchmarks.} Does the two-orders-of-magnitude
improvement survive on H$_2$ / LiH Jordan-Wigner-mapped Hamiltonians or the Schwinger model, whose
Pauli-decompositions are much denser than the nearest-neighbor ZZ+X+Z Ising chain? The paper's
brickwall only has nearest-neighbor ZZ bricks; expressiveness on long-range XXYY / YYYY / ZZZZ terms
is untested.
\emph{Next steps:} Use Qiskit-Nature H$_2$/STO-3G (4 qubits, $\sim 15$ Paulis) and Schwinger $n=4$
($\sim 10$ Paulis) as $U_{\text{target}}$, at $t \in \{0.1, 0.4, 1.0\}$; report whether the $\sim 10^2$
edge holds. Compute: uicgpu A100 (free).

\item \textbf{Noise robustness of the compiled circuit.} Does the classically-optimized brickwall
degrade faster than depth-matched Trotter II under realistic hardware noise (depolarizing per-gate,
coherent $R_{zz}$ over/under-rotation, $T_1/T_2$)? The paper's $\varepsilon_{\text{approx}}$ is a
noiseless metric; fine-tuned optima may be brittle.
\emph{Next steps:} Extend the Qiskit cross-check to \texttt{qiskit-aer} density-matrix simulator with
a \texttt{NoiseModel} at $p_{1q} \in \{10^{-4}, 10^{-3}\}$, $p_{2q} \in \{5{\times}10^{-4}, 5{\times}10^{-3}, 10^{-2}\}$,
plus $R_{zz}$ miscalibration $\delta \in \{10^{-3}, 10^{-2}\}$. Find the noise threshold at which
optimized loses its edge. Compute: local CPU (Aer).

\item \textbf{Amortized / meta-learned optimizer for compilation.} Would a learned optimizer (L2O
MLP or small transformer trained across the $(H, t, L)$ manifold) reduce compilation from
seconds-to-minutes per instance to milliseconds, enabling on-the-fly rather than offline compilation
--- e.g.\ for chemistry sweeps of hundreds of geometries or time-dependent pulse-slicing?
\emph{Next steps:} Generate $\sim 1000$ optimized $(H, t, L, \theta^*)$ triples across random
$(J,g,h)$ and $t \in \{0.05, 0.1, 0.2, 0.4\}$, $n \in \{3, 4\}$. Train a small model to predict
$\theta^*$; measure predicted-$\theta$ $\varepsilon_{\text{approx}}$ vs full-optimizer baseline.
Compute: uicgpu A100 (free) for training.

\end{enumerate}
